\begin{proposition}[5.10.16]
\label{IV.2.s5.10.16-fr}
Sean $X$ un preesquema localmente noetheriano,
\[
X'=\operatorname{Spec}(\mathcal H^0_{X/Z^{(2)}}(\mathcal O_X))
\]
su $Z^{(2)}$-clausura, $g:X'\to X$ el morfismo estructural. Supongamos que $X$
carezca de ciclo primo asociado de codimensión 1.
\begin{enumerate}
\item[\textup{(i)}] Para que, en un punto $x\in X$, $X$ verifique $(S_2)$, es necesario
y suficiente que el morfismo $X'_x\to X_x$ obtenido de $g$ (notaciones de
\textup{(5.10.12)}) sea un isomorfismo. Esta condición se verifica siempre si
$\operatorname{codim}(\overline{\{x\}},X)\leq1$.
\item[\textup{(ii)}] Supongamos además que $g$ sea un morfismo finito \textup{(véanse en
(5.11.2) condiciones} suficientes \textup{para que así sea)}. Entonces el conjunto $U$
de los puntos en los que $X$ verifica $(S_2)$ es abierto y
$\operatorname{codim}(X\setmin U,X)\geq2$; además, $U$ es el mayor abierto de $X$
tal que la restricción $g^{-1}(U)\to U$ de $g$ sea un isomorfismo.
\item[\textup{(iii)}] Bajo las mismas hipótesis que en \textup{(ii)}, $X'$ verifica $(S_2)$ y,
para todo $x'\in X'$ tal que\linebreak
$\operatorname{codim}(\overline{\{x'\}},X')\leq1$, el punto $x=g(x')$ es tal
que
$\operatorname{codim}(\overline{\{x\}},X)
=\operatorname{codim}(\overline{\{x'\}},X')$.
\item[\textup{(iv)}] Siendo las hipótesis las de \textup{(ii)}, sea $\mathcal F$ un
$\mathcal O_X$-módulo coherente de soporte $X$, tal que
$\mathcal H^0_{X/Z^{(2)}}(\mathcal F)$ sea coherente; entonces el
$\mathcal O_{X'}$-módulo $\mathcal F'$ tal que
$g_*(\mathcal F')=\mathcal H^0_{X/Z^{(2)}}(\mathcal F)$ es coherente y
verifica $(S_2)$, y su soporte es una unión de componentes irreducibles de
$X'$.
\end{enumerate}
\end{proposition}

Las afirmaciones (i) y (ii) se incluyen como recordatorio, pues ya han sido
demostradas en esencia: en efecto, (i) resulta de (5.10.12, (i)) y de
(5.10.14), y (ii) es un caso particular de (5.10.15, (ii)).

Demostremos (iii); pongamos $x=g(x')$; como $g$ es finito, también lo es el
morfismo $X'_x\to X_x$, por tanto
$\dim(\mathcal O_{X',x'})\leq\dim(X'_x)\leq\dim(X_x)
=\dim(\mathcal O_{X,x})$
en virtud de (5.4.1). Supongamos primero que
$\dim(\mathcal O_{X',x'})\leq1$ y demostremos que entonces
$\dim(\mathcal O_{X,x})\leq1$. En caso contrario, se tendría
$x\in Z^{(2)}$, por lo que, en virtud de (5.10.12, (ii)) aplicado a
$\mathcal O_X$, se tendría
$\operatorname{prof}(\mathcal O_{X',x'})\geq2$, lo que es absurdo. Se tiene,
por tanto, $x\in X\setmin Z^{(2)}$, y por consiguiente $\mathcal O_{X',x'}$ es
isomorfo a $\mathcal O_{X,x}$ (5.10.12, (i)), de donde
$\dim(\mathcal O_{X',x'})=\dim(\mathcal O_{X,x})$. Además, como $X$ no tiene
ciclos

\oldpage[IV-2]{121}

primos asociados de codimensión 1, la hipótesis
$\dim(\mathcal O_{X,x})=1$ implica $x\notin\operatorname{Ass}(\mathcal O_X)$,
por tanto $\operatorname{prof}(\mathcal O_{X,x})=1$, y por consiguiente
también $\operatorname{prof}(\mathcal O_{X',x'})=1$. Supongamos ahora
$\dim(\mathcal O_{X',x'})\geq2$, por tanto
$\dim(\mathcal O_{X,x})\geq2$, es decir, $x\in Z^{(2)}$; se deduce de ello que
$\operatorname{prof}(\mathcal O_{X',x'})\geq2$ por (5.10.12, (ii)). Esto
establece las afirmaciones de (iii).

Para probar (iv), basta reemplazar $\mathcal O_X$ por $\mathcal F$ en el
razonamiento precedente, que establece que $\mathcal F'$ verifica $(S_2)$ y
que, si $\dim(\mathcal F'_{x'})\leq1$, $\mathcal F_x$ y
$\mathcal F'_{x'}$ son di-isomorfos; en particular, si
$\dim(\mathcal F'_{x'})=0$, se tiene $\dim(\mathcal F_x)=0$, por tanto
$\dim(\mathcal O_{X,x})=0$, puesto que $\mathcal F$ tiene por soporte $X$, y
finalmente $\dim(\mathcal O_{X',x'})=0$; toda componente irreducible de
$\operatorname{Supp}(\mathcal F')$ es, por tanto, una componente irreducible
de $X'$, puesto que $\mathcal F'$ es coherente, y por tanto
$\operatorname{Supp}(\mathcal F')$ es cerrado.

\begin{proposition}[5.10.17]
\label{IV.2.s5.10.17-fr}
Sean $A$ un anillo íntegro noetheriano, y designemos por $A^{(1)}$ la
intersección de los anillos locales $A_{\mathfrak p}$, donde $\mathfrak p$
recorre el conjunto de los ideales primos de $A$ de altura 1. Supongamos que
$A^{(1)}$ sea una $A$-álgebra \emph{finita}. Entonces:
\begin{enumerate}
\item[\textup{(i)}] El anillo $A^{(1)}$ verifica la condición $(S_2)$.
\item[\textup{(ii)}] El conjunto $U$ de los $\mathfrak p\in\operatorname{Spec}(A)$
tales que el homomorfismo canónico
$A_{\mathfrak p}\to(A^{(1)})_{\mathfrak p}$ sea biyectivo es igual al
conjunto de los $\mathfrak p$ tales que $A_{\mathfrak p}$ verifica $(S_2)$;
$U$ es abierto en $X=\operatorname{Spec}(A)$ y se tiene
$\operatorname{codim}(X\setmin U,X)\geq2$.
\item[\textup{(iii)}] Para toda parte multiplicativa $S$ de $A$, $(S^{-1}A)^{(1)}$ es
una $(S^{-1}A)$-álgebra finita.
\item[\textup{(iv)}] Sea $B$ una $A$-álgebra finita, íntegra y que contiene a $A$.
Entonces $B^{(1)}$ es una $B$-álgebra finita. Además, para todo ideal primo
$\mathfrak q$ de $B$, de altura 1, el ideal primo $\mathfrak q\cap A$ de $A$
es de altura 1.
\end{enumerate}
\end{proposition}

Teniendo en cuenta la fórmula (5.9.3.1), se ve que
$X'=\operatorname{Spec}(A^{(1)})$ es la $Z^{(2)}$-\emph{clausura} de
$X=\operatorname{Spec}(A)$; como $A$ no tiene ideales primos asociados
inmersos, las propiedades (i) y (ii) son casos particulares de (5.10.16,
(i), (ii) y (iii)). Para probar (iii), basta observar que se tiene
$(S^{-1}A)^{(1)}=S^{-1}A^{(1)}$, lo que es un caso particular de (5.9.4): en
efecto, $S^{-1}A$ es un $A$-módulo plano, los ideales primos de $S^{-1}A$ son
los ideales $S^{-1}\mathfrak p$, donde
$\mathfrak p\in\operatorname{Spec}(A)$ no interseca $S$, y se tiene
$\operatorname{ht}(S^{-1}\mathfrak p)=\operatorname{ht}(\mathfrak p)$. Como
$A^{(1)}$ es una $A$-álgebra finita, $S^{-1}A^{(1)}$ es una
$S^{-1}A$-álgebra finita, de donde (iii).

Para demostrar (iv), pongamos $Y=\operatorname{Spec}(B)$, y sea
$f:Y\to X$ el morfismo estructural; como es finito, resulta de (5.4.1) que,
para todo $y\in Y$, se tiene
$\dim(\mathcal O_{Y,y})\leq\dim(\mathcal O_{X,f(y)})$; por tanto, si
$T=f^{-1}(Z^{(2)}(X))$, se tiene $T\supset Z^{(2)}(Y)$. Demostremos que
$\mathcal G=\mathcal H^0_{X/Z^{(2)}(X)}(f_*(\mathcal O_Y))$
es coherente; en efecto, $f_*(\mathcal O_Y)=\widetilde B$, considerando $B$
como $A$-módulo; pero, como $B$ es una $A$-álgebra íntegra finita, su cuerpo
de fracciones es finito sobre el cuerpo de fracciones de $A$, por tanto $B$
está contenido en un $A$-módulo libre de tipo finito, y por consiguiente
(5.9.2, (i)) $\mathcal G$ es un sub-$\mathcal O_X$-módulo cuasicoherente de
$(\mathcal H^0_{X/Z^{(2)}(X)}(\mathcal O_X))^n$ para un $n$ adecuado; siendo
este último coherente por hipótesis, ocurre lo mismo con $\mathcal G$. Ahora
bien, resulta de la definición (5.9.1.2) que $\mathcal G$ es isomorfo a
$f_*(\mathcal H^0_{Y/T}(\mathcal O_Y))$; esto prueba \emph{a fortiori} que
$\mathcal H^0_{Y/T}(\mathcal O_Y)$ es un $\mathcal O_Y$-módulo coherente. Se
deduce entonces de (5.10.10, (ii)), aplicado a $\mathcal O_Y$ y al punto
genérico de $Y$, que se tiene $\operatorname{codim}(T,Y)\geq2$, es decir,
$T\subset Z^{(2)}(Y)$, y finalmente $T=Z^{(2)}(Y)$. Esto demuestra las dos
afirmaciones de (iv).

\oldpage[IV-2]{122}
